\subsubsection{CU-03: Ver Sitios de Reparación Cercanos}
Este caso de uso describe el proceso mediante el cual un usuario podrá 
consultar sitios de reparación cercanos a su ubicación actual, asociados 
al tipo de producto escaneado, con el objetivo de fomentar la cultura 
de reparación y la sustentabilidad.


\begin{table}[H]
\centering
\begin{tabular}{p{3.5cm} p{10cm} p{1cm}}\toprule \textbf{Caso de Uso} &
\textit{Visitar Sitios Cercanos de Reparación} & \textit{03} \\

\midrule
\textbf{Actores}       & Usuario             \\[6pt]
\textbf{Tipo}          & Primario            \\[6pt]
\textbf{Precondición}  & El usuario debe haber escaneado previamente un producto
y haber habilitado el permiso de ubicación.  \\[6pt]
\textbf{Postcondición} & El sistema muestra un listado de sitios de reparación
cercanos a la ubicación actual del usuario.  \\[6pt]

\midrule
\multicolumn{3}{l}{\textbf{Propósito}} \\*[2pt]
\multicolumn{3}{p{14.5cm}}{Permitir que el usuario encuentre sitios de
reparación cercanos a su ubicación actual.} \\[6pt]

\midrule
\multicolumn{3}{l}{\textbf{Resumen}} \\*[2pt]
\multicolumn{3}{p{14.5cm}}{El usuario, tras escanear un producto, puede
consultar sitios de reparación cercanos a su ubicación. El sistema obtiene la
posición actual del dispositivo, consulta una API externa y muestra un listado
con nombre, dirección, distancia, teléfono, horario y calificación de cada
establecimiento. El usuario puede seleccionar uno y obtener la ruta directamente
desde la aplicación.} \\[6pt]
\bottomrule
\end{tabular}
\end{table}

\begin{table}[H]
\centering
\begin{tabular}{p{15.5cm}}\toprule \textbf{Descripción} \\
\midrule
\textbf{Flujo Principal} \\[2pt]
\begin{enumerate}
  \item El usuario escanea el producto y accede a su información.
  \item El usuario selecciona la opción \textit{Ver sitios de reparación
      cercanos}.
  \item El sistema obtiene la ubicación actual del dispositivo del usuario.
  \item El sistema consulta la API de sitios de reparación con la ubicación y
      la categoría del producto escaneado.
  \item El sistema muestra un listado de sitios de reparación con la siguiente
      información por cada uno:
    \begin{enumerate}
        \item \textbf{Nombre} del establecimiento.
        \item \textbf{Dirección} y distancia aproximada al usuario.
        \item \textbf{Teléfono} de contacto.
        \item \textbf{Horario} de atención.
        \item \textbf{Calificación} promedio del establecimiento.
    \end{enumerate}
  \item El usuario selecciona un sitio de reparación de la lista.
  \item El sistema abre el mapa del dispositivo con la ruta hacia el sitio
      seleccionado.
\end{enumerate} \\
\midrule
\textbf{Flujo Alterno} \\[2pt]
\begin{enumerate}
\item[3a.] Si el usuario no ha habilitado el permiso de ubicación, el
    sistema muestra un mensaje solicitando activarlo y redirige a la
    configuración del dispositivo.
\item[4a.] Si no se encuentran sitios de reparación cercanos, el sistema
    muestra un mensaje informativo.
\item[4b.] Si no hay conexión a internet, el sistema muestra un mensaje de
    error indicando que se requiere conexión para realizar la búsqueda.
\end{enumerate}
\\
\bottomrule
\end{tabular}
\end{table}